% ===== Block I: PDF p.38–42 | 原文 §4 Further directions（4.1 Moduli of hyperbolic surfaces；4.2 String theory and moduli of maps） | 唯一编号公式 4.1–4.14 =====
% 编号环境：Theorem 4.1、Theorem 4.2、Exercise 4.1、Theorem 4.3；图 15 略（仅图注中译）；表 3 照原文抄录重排
% [NOTE: 本块从 §4 与 §4.1 的标题之后开始（这两个标题由 Lead 写在 main.tex 中）；文件末尾自带 §4.2 的标题（全项目唯一例外）]
% [原文 p.38]

在 JT 引力中，理论的路径积分是在双曲度规的空间（而非复结构的空间）上进行的。换言之，“正确”的模空间是\emph{双曲结构}的模空间：
\begin{equation}
  \cM^{\mathrm{hyp}}_{g,n}(L_1,\dots,L_n)=
  \left\{ X \;\middle|\;
  \begin{array}{l}
    \mbox{$X$ 是亏格为 $g$ 的双曲曲面} \\[2pt]
    \mbox{带有 $n$ 条标号的测地边界} \\[2pt]
    \mbox{其长度为 $L_1,\dots,L_n$}
  \end{array}
  \right\}\Big/\sim
  \tag{4.1}
\end{equation}
其中 $X\sim X'$ 当且仅当存在从 $X$ 到 $X'$ 的等距，且保持边界分支的标号。
% [Tutor Note: (4.1) 的商只对“保持边界分支标号”的等距取，故标号是数据的一部分，模空间确实依赖有序的 $L=(L_1,\dots,L_n)$。核对维数：$2(3g-3+n)$ 个实维恰为 $\cM_{g,n}$ 的复维 $3g-3+n$ 的两倍；两个极端核对：$(g,n)=(0,3)$ 给出 0 维（给定三条测地边界长度的双曲裤子在等距下唯一），$(g,n)=(1,1)$ 给出 2 维（固定边界长度的单孔环面模空间）。另：原文条件行“with $n$ labelled geodesics boundaries”中“geodesics”为笔误，应为“geodesic”（原文笔误，见 ERRATA）。]

这与黎曼曲面的模空间有什么关系？一个非平凡的结果——它是 Riemann 单值化定理的推论——是 $\cM^{\mathrm{hyp}}_{g,n}(L)$ 同胚于第 2 节所讨论的黎曼曲面模空间。

\begin{theorem}[4.1]
空间 $\cM^{\mathrm{hyp}}_{g,n}(L)$ 是一个维数为 $2(3g-3+n)$ 的光滑实轨形（smooth real orbifold）。此外，对一切 $L\in\R_+^n$，它（作为光滑实轨形）同胚于光滑黎曼曲面的模空间：
\begin{equation}
  \cM^{\mathrm{hyp}}_{g,n}(L)\cong\cM_{g,n}.
  \tag{4.2}
\end{equation}
\end{theorem}
% [Tutor Note: “对一切 $L\in\R_+^n$ 同胚”是强陈述：不同的 $L$ 只改变辛形式而不改变光滑结构，这正是下一段能定义 $\VWP_{g,n}(L)$ 的前提。核对 $g=0,n=3$：$\cM_{0,3}$ 为单点，而给定三条边界长度的双曲裤子唯一，与 (4.2) 相符；$g=1,n=1$ 时两边都是 2 维实流形。]

对任意固定的 $L\in\R_+^n$，模空间 $\cM^{\mathrm{hyp}}_{g,n}(L)$ 自然地配备一个辛形式，称为 Weil--Petersson 形式，记作 $\omega_{\WP}$。特别地，我们可以定义体积
\begin{equation}
  \VWP_{g,n}(L)=\int_{\cM^{\mathrm{hyp}}_{g,n}(L)}\frac{\omega_{\WP}^{3g-3+n}}{(3g-3+n)!}.
  \tag{4.3}
\end{equation}

这类结构的一个玩具例子是：用半径 $L$ 的球面 $S^2(L)$ 对 $\R_+\ni L$ 作纤维化。尽管所有纤维都同胚于 $\PP^1$，每个纤维却带有依赖于底空间上点 $L$ 的特定辛几何。例如，$S^2(L)$ 的面积为 $4\pi L^2$。然而，我们可以把这特定的几何搬到 $\PP^1$ 上，得到 $\PP^1$ 上一个依赖于 $L$ 的表达式。例如，在球极投影给出的同构 $S^2(L)\cong\PP^1$ 之下，我们发现 $S^2(L)$ 上的辛形式被映为如下关于 $L$ 的多项式：
\begin{equation}
  4L^2\,\Re\frac{\dd z\,\dd\bar z}{(1+|z|^2)^2}.
  \tag{4.4}
\end{equation}
% [Tutor Note: 核对 (4.4) 的归一化：令 $z=x+\ii y$，把 $\dd z\,\dd\bar z$ 理解为对称的 $|\dd z|^2=\dd x^2+\dd y^2$（若按外积 $\dd z\wedge\dd\bar z$ 理解则实部为零），则 $\omega=4L^2\,\dd x\,\dd y/(1+r^2)^2$，而 $\int_{\R^2}\dd x\,\dd y/(1+r^2)^2=\pi$，故 $\int_{S^2(L)}\omega=4\pi L^2$，与半径 $L$ 的球面面积一致。]

对 Weil--Petersson 形式与同构 $\cM^{\mathrm{hyp}}_{g,n}(L)\cong\cM_{g,n}$ 而言，相应的结果归功于 Wolpert（$L_i=0$ 的情形）与 Mirzakhani（一般情形）[Wol85; Mir07b]。

\begin{theorem}[4.2]
在同胚 $\cM^{\mathrm{hyp}}_{g,n}(L)\cong\cM_{g,n}$ 之下，Weil--Petersson 形式作为闭形式延拓到 $\Mc_{g,n}$，并定义上同调类
\begin{equation}
  2\pi^2\kappa_1+\frac12\sum_{i=1}^n L_i^2\psi_i.
  \tag{4.5}
\end{equation}
\end{theorem}
% [Tutor Note: (4.5) 是 $\Mc_{g,n}$ 上的 2-上同调类（$\kappa_1$ 与 $\psi_i$ 都是 2 次类）。核对 $g=n=1$：$\Mc_{1,1}$ 上 $\int\psi_1=\frac1{24}$，由 (4.5) 得 $\partial V/\partial L=\int L\psi_1=\frac{L}{24}$，而表中 $\VWP_{1,1}(L)=\frac{L^2}{48}+\frac{\pi^2}{12}$ 对 $L$ 求导恰为 $\frac{L}{24}$，两者自洽。]

% [原文 p.39]
\figomit{15}{测地线 $\alpha_p$（红色）及其若干可能的行为，连同它所确定的简单闭曲线（绿色）。左图：弧 $\alpha_p$ 与边界分支 $\partial_m\Sigma$ 相交（$\mathrm{b}_m$ 型），并确定一条简单闭曲线 $\gamma$。另外两种情形中，$\alpha_p$ 分别与 $\partial_1\Sigma$ 及自身相交（$\mathrm{c}$ 型），确定两条简单闭曲线 $(\gamma,\gamma')$。}
% [Tutor Note: 图注中“$\alpha_p$ intersect $\partial_1\Sigma$ and itself respectively”英文语法应为“intersects”（原文笔误，译文照译）；正文称 $\alpha_p$ 为 geodesic，图注称 arc，指同一对象。$\mathrm{b}_m$ 型与 $\mathrm{c}$ 型的区别是：前者 $\alpha_p$ 撞上另一条边界分支 $\partial_m\Sigma$，后者 $\alpha_p$ 撞上 $\partial_1\Sigma$ 或自身，后者切出两条简单闭曲线。]

上文结果的一个直接推论是：Weil--Petersson 体积是有限的（这一点并不显然，因为 $\cM^{\mathrm{hyp}}_{g,n}(L)$ 不紧），并且是边界长度平方的对称多项式，其系数是涉及 $\psi$ 类与 $\exp(2\pi^2\kappa_1)$ 的相交数：
\begin{equation}
  \VWP_{g,n}(L)=\sum_{\substack{d_1,\dots,d_n\ge0\\ d_1+\cdots+d_n\le 3g-3+n}}\int_{\cM_{g,n}}\ee^{2\pi^2\kappa_1}\prod_{i=1}^n \psi_i^{d_i}\,\frac{L_i^{2d_i}}{2^{d_i}d_i!}.
  \tag{4.6}
\end{equation}
% [Tutor Note: 体积有限性非平凡：固定边界长度时 $\cM^{\mathrm{hyp}}_{g,n}(L)$ 不紧（曲面可沿简单闭测地线无限“捏瘦”退化），而 Weil--Petersson 体积仍有限，这正是 Mirzakhani 工作的出发点之一。]
% [Tutor Note: (4.6) 可由 (4.5) 直接得到：在 $\cM_{g,n}$ 上积分 $\exp\big(2\pi^2\kappa_1+\frac12\sum_i L_i^2\psi_i\big)$，按 $L_i$ 展开，$L_i^{2d_i}$ 的系数正是 $\frac{1}{2^{d_i}d_i!}\psi_i^{d_i}$。这也解释了求和截断 $d_1+\cdots+d_n\le 3g-3+n$：$\dim_\C\cM_{g,n}=3g-3+n$，更高次的 $\psi$ 类乘积积分为零。]

这些相交数恰好具有 CohFT 关联函数的形式，因而可以借助拓扑递归计算！

\begin{exercise}[4.1]
考虑谱曲线
\begin{equation}
  \left(\PP^1,\ x(z)=\frac{z^2}{2},\ y(z)=\frac{\sin(2\pi z)}{2\pi},\ \omega_{0,2}(z_1,z_2)=\frac{\dd z_1\dd z_2}{(z_1-z_2)^2}\right),
  \tag{4.7}
\end{equation}
另见习题 3.4。证明：与上述谱曲线相应的拓扑递归关联函数计算 Weil--Petersson 体积的 Laplace 变换的微分：
\begin{equation}
  \omega_{g,n}(z_1,\dots,z_n)=\dd z_1\cdots\dd z_n\left(\prod_{i=1}^n\int_0^\infty \dd L_i\,\ee^{-z_iL_i}\right)\VWP_{g,n}(L_1,\dots,L_n).
  \tag{4.8}
\end{equation}
\end{exercise}
% [Tutor Note: 这里把谱曲线连同 $\omega_{0,2}$ 一起给出：$\omega_{0,2}$ 本身是二阶极点的 $\dd$-微分，并非 $y$ 的导数；$x=z^2/2$ 只是 $x$-坐标的标定（$\dd x=z\,\dd z$），而 $y(z)=\frac{\sin(2\pi z)}{2\pi}$ 在 $z=0$ 附近 $y\sim z$，故 $y$ 在 $x$ 的平方根分支上是单值的。本习题即 Mirzakhani 递归的 TR 形式，见定理 4.3 后的说明与 [EO]。]
% [Tutor Note: 核对 (4.8) 的一个小例子：$g=0,n=3$ 时 $\VWP_{0,3}=1$，于是 (4.8) 给出 $\omega_{0,3}=\dd z_1\dd z_2\dd z_3/(z_1z_2z_3)$（每个 Laplace 积分 $\int_0^\infty\dd L\,\ee^{-zL}=1/z$）；这是对 TR 输出的一条可检验预测，也说明式中 $\dd z_1\cdots\dd z_n$ 与“微分”一词的对应。]

与拓扑递归等价的一个陈述（针对体积本身，而非其 Laplace 变换）由 M. Mirzakhani 在一系列出色的论文中证明 [Mir07a; Mir07b]。她的方法完全是几何的（而非代数的），基于 McShane [McS98] 的如下简单想法。

考虑一个固定的、带测地边界的双曲曲面 $(\Sigma,h)$。相对于双曲测度随机地选取一点 $p\in\partial_1\Sigma$，并考虑从 $p$ 出发、与 $\partial_1\Sigma$ 正交的测地线 $\alpha_p$。那么以下互斥的情形之一必定发生（参看图 15）。
\begin{enumerate}[label=\roman*),leftmargin=2.6em]
  \item 测地线 $\alpha_p$ 与某个 $m\in\{2,\dots,n\}$ 的 $\partial_m\Sigma$ 相交，而不自交。
  \item 测地线 $\alpha_p$ 与 $\partial_1\Sigma$ 相交，而不自交。
% [原文 p.40]
  \item 测地线 $\alpha_p$ 自交。
  \item 测地线 $\alpha_p$ 既不自交，也不与任何边界分支相交（它无限地螺旋下去）。
\end{enumerate}

一方面，由 Birman--Series 的一个结果，情形 (iv) 出现的概率为零。于是我们简单地有 $1=\mathbb{P}_{\mathrm{I}}+\mathbb{P}_{\mathrm{II}}+\mathbb{P}_{\mathrm{III}}$，即总概率表示为出现 (i)、(ii) 或 (iii) 的概率之和。为了计算这些概率，Mirzakhani 如下进行。考虑 $\partial_1\Sigma$、测地线 $\alpha_p$ 以及（仅对情形 (i)）$\partial_m\Sigma$ 的并。这个曲线并的一个充分小的邻域在拓扑上是一对裤子（pair of pants）。通过取边界分支的测地代表，我们得到一对嵌入的双曲裤子，其测地边界为 $(\partial_1\Sigma,\partial_m\Sigma,\gamma)$——我们称之为 $\mathrm{b}_m$ 情形——或为 $(\partial_1\Sigma,\gamma,\gamma')$——我们称之为 $\mathrm{c}$ 情形（再参看图 15）。于是我们可以写
\begin{equation}
  1=\sum_{m=2}^n\mathbb{P}_{\mathrm{B}_m}+\mathbb{P}_{\mathrm{C}}.
  \tag{4.9}
\end{equation}

Mirzakhani 将概率 $\mathbb{P}_{\mathrm{B}_m}$ 与 $\mathbb{P}_{\mathrm{C}}$ 计算为相应裤子的双曲长度的函数，得到了她著名的恒等式：
\begin{equation}
  1=\sum_{m=2}^n\sum_\gamma B(L_1,L_m,\ell(\gamma))+\frac12\sum_{\gamma,\gamma'}C(L_1,\ell(\gamma),\ell(\gamma')),
  \tag{4.10}
\end{equation}
其中 $B$ 与 $C$ 是显式的双曲函数
\begin{align}
  B(L,L',\ell)&=1-\frac1L\log\left(\frac{\cosh(\frac{L'}2)+\cosh(\frac{L+\ell}2)}{\cosh(\frac{L'}2)+\cosh(\frac{L-\ell}2)}\right),\nonumber\\
  C(L,\ell,\ell')&=\frac2L\log\left(\frac{\ee^{\frac L2}+\ee^{\frac{\ell+\ell'}2}}{\ee^{-\frac L2}+\ee^{\frac{\ell+\ell'}2}}\right).
  \tag{4.11}
\end{align}
% [Tutor Note: (4.9) 中 $\mathbb{P}_{\mathrm{B}_m}$ 的下标 $m$ 是被 $\alpha_p$ 撞到的边界分支编号，与 (4.10) 中 $B(L_1,L_m,\ell)$ 的第二个变量 $L_m$ 对应；$B$ 型对应裤子边界 $(\partial_1\Sigma,\partial_m\Sigma,\gamma)$，$C$ 型对应 $(\partial_1\Sigma,\gamma,\gamma')$。$(\gamma,\gamma')$ 无序，故 (4.10) 第二项前有因子 $1/2$，否则同一对曲线会被计数两次。]
% [Tutor Note: 用退化极限核对 $B,C$：$\ell=0$ 时 $B(L,L',0)=1-\frac1L\log 1=1$；$\ell=\ell'=0$ 时 $C(L,0,0)=\frac2L\log\frac{\ee^{L/2}+1}{\ee^{-L/2}+1}=\frac2L\log\ee^{L/2}=1$，与 (4.10) 左端为 $1$、各项为概率的解释自洽。]

将常数函数 1 在模空间上积分，左边给出 Weil--Petersson 体积；而右边由于去掉了裤子，可以表示为涉及更低拓扑复杂度的体积的特定积分公式。

\begin{theorem}[4.3（Mirzakhani 递归）]
Weil--Petersson 体积由以下关于 $2g-2+n>1$ 的递归唯一确定
\begin{align}
  \VWP_{g,n}(L_1,\dots,L_n)&=\sum_{m=2}^n\int_0^\infty \dd\ell\,\ell\,B(L_1,L_m,\ell)\,\VWP_{g,n-1}(\ell,L_2\dots,\widehat{L_m},\dots,L_n)\nonumber\\
  &\quad+\frac12\int_0^\infty\!\!\int_0^\infty \dd\ell\,\dd\ell'\,\ell\ell'\,C(L_1,\ell,\ell')\Bigg(\VWP_{g-1,n+2}(\ell,\ell',L_2,\dots,L_n)\nonumber\\
  &\qquad+\sum_{\substack{g_1+g_2=g\\ I_1\sqcup I_2=\{2,\dots,n\}}}\VWP_{g_1,1+|I_1|}(\ell,L_{I_1})\,\VWP_{g_2,1+|I_2|}(\ell',L_{I_2})\Bigg)
  \tag{4.12}
\end{align}
并采用约定 $\VWP_{0,1}=\VWP_{0,2}=0$ 以及初值 $\VWP_{0,3}(L_1,L_2,L_3)=1$、$\VWP_{1,1}(L)=\frac{L^2}{48}+\frac{\pi^2}{12}$。
\end{theorem}
% [Tutor Note: 递归只对 $2g-2+n>1$（即 $3g-3+n>0$）陈述：此时模空间非空且维数为正。约定 $\VWP_{0,1}=\VWP_{0,2}=0$ 是因为 $2g-2+n\le1$ 的双曲曲面不存在（空模空间），而 $g=0,n=3$ 与 $g=1,n=1$ 是递归的起点。]
% [Tutor Note: (4.12) 三项的几何来源是沿 (4.10) 中的曲线把裤子粘回去的三种方式。用边界条数与 Euler 示性数核对：第一项把 $\partial_1,\partial_m$ 换成新边界 $\ell$，边界数 $n\to n-1$，故出现 $\VWP_{g,n-1}$（比左边少 1 恰因移除了一条裤子，$\chi=-1$）；第二项对应沿一条非分离简单闭测地线切开，亏格 $g\to g-1$、边界数 $n\to n+2$，此时 $2g-2+n$ 保持不变；第三项按 $g_1+g_2=g$、$I_1\sqcup I_2=\{2,\dots,n\}$ 把曲面分成两块，$L_{I_1},L_{I_2}$ 分别承载两块上的其余边界长度。]

式 (4.7) 给出的谱曲线上的拓扑递归恰好产生 Mirzakhani 递归的 Laplace 变换 [EO]。

% [原文 p.41]
\clearpage
\begin{table}[t]\centering\small
\caption*{表 3：用拓扑递归计算出的 Weil--Petersson 多项式 $\VWP_{g,n}(L)$ 一览。这里 $m_\lambda$ 是与分拆 $\lambda$ 相联系的单项对称多项式，取值于 $L_1^2,\dots,L_n^2$。}
\begin{tabular}{@{}c|>{\raggedright\arraybackslash}p{\dimexpr\textwidth-2.9cm\relax}@{}}
\toprule
$(g,n)$ & $\VWP_{g,n}(L_1,\dots,L_n)$ \\
\midrule
$(0,3)$ & $1$ \\
$(0,4)$ & $\frac12m_{(1)}+2\pi^2$ \\
$(0,5)$ & $\frac18m_{(2)}+\frac12m_{(1^2)}+3\pi^2m_{(1)}+10\pi^4$ \\
$(0,6)$ & $\frac1{48}m_{(3)}+\frac3{16}m_{(2,1)}+\frac34m_{(1^3)}+\frac{3\pi^2}{2}m_{(2)}+6\pi^2m_{(1^2)}+26\pi^4m_{(1)}+\frac{244\pi^6}{3}$ \\
$(0,7)$ & $\frac1{384}m_{(4)}+\frac1{24}m_{(3,1)}+\frac3{32}m_{(2^2)}+\frac38m_{(2,1^2)}+\frac32m_{(1^4)}+\frac{5\pi^2}{12}m_{(3)}+\frac{15\pi^2}{12}m_{(2,1)}$ \newline
$\quad+15\pi^2m_{(1^3)}+20\pi^4m_{(2)}+80\pi^4m_{(1^2)}+\frac{910\pi^6}{3}m_{(1)}+\frac{2758\pi^8}{3}$ \\
\midrule
$(1,1)$ & $\frac1{48}m_{(1)}+\frac{\pi^2}{12}$ \\
$(1,2)$ & $\frac1{192}m_{(2)}+\frac1{96}m_{(1^2)}+\frac{\pi^2}{12}m_{(1)}+\frac{\pi^4}{4}$ \\
$(1,3)$ & $\frac1{1152}m_{(3)}+\frac1{192}m_{(2,1)}+\frac1{96}m_{(1^3)}+\frac{\pi^2}{24}m_{(2)}+\frac{\pi^2}{8}m_{(1^2)}+\frac{13\pi^4}{24}m_{(1)}+\frac{14\pi^6}{9}$ \\
$(1,4)$ & $\frac1{9216}m_{(4)}+\frac1{768}m_{(3,1)}+\frac1{384}m_{(2^2)}+\frac1{128}m_{(2,1^2)}+\frac1{64}m_{(1^4)}+\frac{7\pi^2}{576}m_{(3)}$ \newline
$\quad+\frac{\pi^2}{12}m_{(2,1)}+\frac{\pi^2}{4}m_{(1^3)}+\frac{41\pi^4}{96}m_{(2)}+\frac{17\pi^4}{12}m_{(1^2)}+\frac{187\pi^6}{36}m_{(1)}+\frac{529\pi^8}{36}$ \\
\midrule
$(2,1)$ & $\frac1{442368}m_{(4)}+\frac{29\pi^2}{138240}m_{(3)}+\frac{139\pi^4}{23040}m_{(2)}+\frac{169\pi^6}{2880}m_{(1)}+\frac{29\pi^8}{192}$ \\
$(2,2)$ & $\frac1{4423680}m_{(5)}+\frac1{294912}m_{(4,1)}+\frac{29}{2211840}m_{(3,2)}+\frac{11\pi^2}{276480}m_{(4)}+\frac{29\pi^2}{69120}m_{(3,1)}+\frac{7\pi^2}{7680}m_{(2^2)}$ \newline
$\quad+\frac{19\pi^4}{7680}m_{(3)}+\frac{181\pi^4}{11520}m_{(2,1)}+\frac{551\pi^6}{8640}m_{(2)}+\frac{7\pi^6}{36}m_{(1^2)}+\frac{1085\pi^8}{1728}m_{(1)}+\frac{787\pi^{10}}{480}$ \\
\midrule
$(3,1)$ & $\frac1{53508833280}m_{(7)}+\frac{77\pi^2}{9555148800}m_{(6)}+\frac{3781\pi^4}{2786918400}m_{(5)}+\frac{47209\pi^6}{418037760}m_{(4)}+\frac{127189\pi^8}{26127360}m_{(3)}$ \newline
$\quad+\frac{8983379\pi^{10}}{87091200}m_{(2)}+\frac{8497697\pi^{12}}{9331200}m_{(1)}+\frac{9292841\pi^{14}}{4082400}$ \\
\bottomrule
\end{tabular}
\end{table}
% [Tutor Note: 逐行核对表格与定理 4.3 的初值：(0,3) 行为 $1$；$n=1$ 时 $m_{(1)}=L_1^2$，(1,1) 行给出 $\frac{1}{48}m_{(1)}+\frac{\pi^2}{12}=\frac{L^2}{48}+\frac{\pi^2}{12}$，与定理 4.3 的初值一致；(0,4) 行为 $\frac12(L_1^2+\cdots+L_4^2)+2\pi^2$，其常数项 $2\pi^2$ 即已知的 $V_{0,4}(0,\dots,0)$。表中 $m_{(1^2)}=m_{(1,1)}=\sum_{i<j}L_i^2L_j^2$，$m_{(2)}=\sum_iL_i^4$，即单项对称多项式在 $L_1^2,\dots,L_n^2$ 处取值。]

\subsection{弦论与映射的模空间（String theory and moduli of maps）}

正如正文中所提到的，拓扑弦论也与黎曼曲面的模空间密切相关。拓扑弦论（用数学术语说，即 Gromov--Witten 理论）旨在计算固定目标时空 $X$ 中弦的世界面——即作为参数化黎曼曲面的世界面，也就是映射
\begin{equation}
  f\colon(\Sigma,p_1,\dots,p_n)\longrightarrow X.
  \tag{4.13}
\end{equation}
这里 $p_1,\dots,p_n$ 是 $\Sigma$ 上的标记点，可以看作世界面 $\Sigma$ 的初始/末态。理论的路径积分便是对这些映射的模空间上的积分：
\begin{equation}
  \cM_{g,n}(X,\beta)=\left\{(\Sigma,p_1,\dots,p_n,f)\;\middle|\;
  \begin{array}{l}
    f\colon(\Sigma,p_1,\dots,p_n)\to X\\
    f_*[\Sigma]=\beta
  \end{array}
  \right\}\Big/\sim,
  \tag{4.14}
\end{equation}
其中 $\beta\in H_2(X,\Z)$ 是一个固定的类（称为次数）。$\cM_{g,n}(X,\beta)$ 及其紧化的正确定义是一个极为精细的数学问题（比黎曼曲面的模空间复杂得多）。相应关联函数的计算则更为微妙。
% [Tutor Note: 原文 “As mentioned in the text” 指引言中“弦的世界面即黎曼曲面”的论述；数学上 $\cM_{g,n}(X,\beta)$ 通常要求 $X$ 为光滑射影簇（或辛流形），否则 Gromov--Witten 不变量的定义需要额外的技术条件（如虚拟基本类）。]
% [Tutor Note: $X=\{*\}$ 时 $\beta=0$，且 $\cM_{g,n}(\{*\},0)=\Mc_{g,n}$，故 Witten 猜想确为 $X$ 为一点的特例，与“冰山一角”的说法对应。(4.14) 中 $f_*[\Sigma]\in H_2(X,\Z)$ 是基本类的前推（原文称其为 degree）；由于对带标记映射的同构取商，$\cM_{g,n}(X,\beta)$ 是模栈（Deligne--Mumford 型），其紧化即稳定映射模空间，这正是原文所说“proper definition”高度精细的原因。]

Witten 猜想可以看作这一理论的冰山一角：它对应于 $X=\{*\}$ 的情形，即零维目标。Eguchi、Hori 与 Xiong [EHX97] 把 Virasoro 约束推广到点目标，并猜想每个目标的配分函数都服从 Virasoro 条件。在一系列出色的论文中 [OP06a; OP06b; OP06c]，Okounkov 与 Pandharipande 对一维情形给出了完整解，证明了 Eguchi--Hori--Xiong 的猜想。除点与复曲线情形之外，Virasoro 约束也被证明对若干特殊类别的目标（维数任意）成立，即：
% [原文 p.42]（p.41 末句 “…extended the Virasoro” 在本页接续为 “constraints for the point …”，该跨页段落至此结束；为不打断段落，本页标记置于该段落之后）
\begin{itemize}
  \item 对 toric Fano 流形以及满足半单性假设的流形，如 Givental--Teleman 所示 [Giv01a; Tel12]，
  \item 对 toric Calabi--Yau 三维流形则更为显式，这沿循 Bouchard--Klemm--Mari\~no--Pasquetti 的“remodelling 猜想”（remodelling conjecture）[BKMP09]，该猜想现已成为定理 [EO15; FLZ20]。
\end{itemize}
% [Tutor Note: 两个条目中 [Giv01a; Tel12] 分别对应 Givental 与 Teleman 的证明，[EO15; FLZ20] 分别对应 Eynard--Orantin 与 Fang--Liu--Zong 对 BKMP 猜想的证明。另注意原文此处拼写为英式 “remodelling”，而参考文献 [BKMP09] 的标题用美式 “Remodeling”，属原文内部拼写不一致，译文分别照原文保留。]
